\chapter{Bayesian Evidence as Counting}
\label{ch:bayes}

Vertical likelihood Monte Carlo was invented for a Bayesian problem, the computation of the marginal likelihood or evidence $Z = \int L(\theta) \, \pi(\theta) \, d\theta$, and only later recognised as a counting method. This chapter closes the loop. It shows that the evidence is a counting problem, that the standard estimators of the evidence are instances of the vertical walk, and that the complexity theory of Part III applies to Bayesian computation with no change, so that the question whether a posterior can be sampled in polynomial time is the same question as whether a $\sharpP$ function can be approximated.

\section{The evidence}

Let $\pi$ be a prior on a parameter space $\Theta$ and $L(\theta) = p(y \mid \theta)$ the likelihood of data $y$. The evidence $Z = \int L \, d\pi$ is the normalising constant of the posterior $\pi(\theta \mid y) = L(\theta) \pi(\theta) / Z$, and ratios of evidences across models are Bayes factors. The evidence is an expectation of the likelihood under the prior, which is the vertical identity with $P = \pi$:
\[
Z = \EE_\pi[L(\theta)] = \int_0^1 \Lambda(s) \, ds ,
\]
where $\Lambda(s)$ is the $(1 - s)$ quantile of $L(\theta)$ under the prior. If $\Theta$ is discrete and $L$ is an indicator, $Z$ is a count divided by $|\Theta|$; if $L$ is a Boltzmann\index{Boltzmann distribution} weight $e^{-S/\gamma}$, $Z$ is a partition function. There is no formal difference between evidence and counting, and the staircase of the evidence is the staircase of the count.

\section{Estimators of the evidence}

The literature on evidence estimation is large, and the vertical view organises it. The harmonic mean estimator of Newton and Raftery, $\hat Z = ( n^{-1} \sum_i 1 / L(\theta_i) )^{-1}$ with $\theta_i$ from the posterior, has infinite variance in general because it weights by the reciprocal likelihood, which is the quantile function\index{quantile function} inverted; in vertical terms it tries to estimate the integral of $\Lambda$ from draws concentrated at $s = 0$, where $\Lambda$ is large, and knows nothing about the width of the region. Importance sampling from a density $g$ estimates $Z$ as the mean of $L \pi / g$ and is the vertical importance sampler with $\Omega$ induced by $g$. Bridge sampling of Meng and Wong estimates the ratio of two normalising constants from samples of both, which is one rung of the ladder. Thermodynamic integration, or path sampling in the sense of Gelman and Meng, writes
\[
\log Z = \int_0^1 \EE_{\pi_t}\big[ \log L(\theta) \big] \, dt , \qquad \pi_t \propto L^t \pi ,
\]
which is the Boltzmann ladder of Chapter~\ref{ch:product} integrated in the inverse temperature coordinate; annealed importance sampling of Neal runs the same ladder sequentially with importance weights; and nested sampling\index{nested sampling} of Skilling\index{Skilling, John}~\cite{Skilling} runs the ladder in the quantile coordinate directly. Polson and Scott's~\cite{PolsonScott} observation is that all of these are choices of $\Omega$, and that the geometric grid of nested sampling is a particularly good one when $\Lambda$ falls off exponentially, as it does for a likelihood that concentrates.

\section{Posterior sampling as level set sampling}

The conditioned draws of the vertical walk, from $\pi$ restricted to $\{ L > u \}$, are draws from the prior truncated to a likelihood level set. This is the auxiliary variable construction of slice sampling\index{slice sampling}, in which the posterior is the marginal of the joint $\pi(\theta) \Ind\{ u < L(\theta) \}$ on $(\theta, u)$, and Damien, Wakefield and Walker's~\cite{DWW} Gibbs sampler alternates $u \mid \theta$, uniform on $(0, L(\theta))$, with $\theta \mid u$, the prior on the level set. Nested sampling is slice sampling with the slice variable driven deterministically down the geometric grid rather than resampled; the two are the same chain on the same joint with different schedules for $u$.

The complexity of the level set draw is the complexity of the posterior. For a log concave likelihood and prior the level sets are convex and the random walk of Applegate\index{Applegate, Kannan and Polson}, Kannan and Polson~\cite{AKP}, and the Metropolis\index{Metropolis algorithm} ball walk of Frieze, Kannan and Polson~\cite{FKP}, mix in polynomial time, so the evidence can be approximated in polynomial time by the vertical walk, which is the Bayesian form of the Dyer Frieze Kannan theorem~\cite{DFK}. For a multimodal likelihood the level sets disconnect at some threshold $u$, the walk must jump between modes, and the problem is as hard as the modes are isolated. The posterior of a mixture model with $k$ components has $k!$ symmetric modes separated by regions of low likelihood; the posterior of a neural network has an exponential number of modes; and the posterior of a Boolean model, with likelihood an indicator of consistency with data, is, under a uniform prior, the uniform distribution on the solution set of a constraint satisfaction problem and is exactly the counting problem of Part III.

\section{The Bayesian form of the thesis}

The thesis of Chapter~\ref{ch:prp} translates word for word. The evidence of any model is the integral of a staircase on the unit interval; the staircase has one step per distinct likelihood value, which for a continuous likelihood is a continuum but which the walk discretises to polynomially many; and the whole difficulty of computing the evidence is the difficulty of drawing from the prior restricted to a likelihood level set. If those draws are polynomial for every level, and the ladder has polynomially many rungs with ratios bounded below by an inverse polynomial, then the evidence has an FPRAS\index{FPRAS}. If the likelihood encodes an $\NP$ complete constraint then the draws at the lowest levels are, unless $\NP = \RP$, not polynomial.

The Bayesian reader may find this reassuring or alarming. It is reassuring in that the familiar difficulties of posterior computation, multimodality and the failure of local samplers to mix across modes, are the same difficulties that the complexity theorists have classified, and the classification applies: a posterior is easy to sample if its level sets are connected with good conductance\index{conductance} and hard for local samplers if they shatter, and in the latter case local samplers fail, while for likelihoods that encode $\NP$ complete constraints no polynomial sampler exists unless $\NP = \RP$. It is alarming in that the Bayesian computational project, taken to its limit, includes the computation of evidences for models whose likelihoods are $\NP$ complete constraints, and for those the project is, unless $\NP = \RP$, intractable in the worst case. The honest statement is that Bayesian computation is tractable exactly as far as the level sets of the likelihood are tractable, and that the vertical coordinate is where to look to see whether they are.

\section{The horseshoe and sparse posteriors}

A case of interest to one of us is the horseshoe posterior for a sparse regression, where the likelihood is Gauss\index{Gauss, Carl Friedrich}ian and the prior is a scale mixture that puts most of its mass near zero with heavy tails. The posterior is not log concave, since the prior is not, but the Gibbs sampler with slice augmentation of the scale variables is widely used in practice, and we claim no mixing result for it. We offer the rest as a heuristic and do not prove it. The level sets of the likelihood are ellipsoids and the level sets of the prior are star shaped regions around the coordinate axes; we expect, but have not shown, that the posterior level sets are connected when the signal is strong. On that expectation the vertical walk for the evidence of a horseshoe regression, which gives the Bayes factor for the inclusion of a set of variables, would be polynomial, and the staircase it climbs would be expected to be smooth, with the drop from the prior bulk to the posterior mode governed by the signal to noise ratio. On that expectation this would be the easy end of the Bayesian spectrum, easy for the reason that the convex body problem is easy, namely that the geometry of the level sets is benign, which we have not shown.

\section{Model selection as counting}

Finally, the choice between models is itself a counting problem when the model space is combinatorial. Variable selection among $p$ predictors has $2^p$ models, and the posterior probability of a model is its evidence times its prior divided by the sum over all models, which is a partition function on $\{0,1\}^p$. Computing the marginal inclusion probability of a variable is computing a ratio of two such partition functions, and for a Gaussian linear model with conjugate priors each evidence is in closed form, so the problem reduces to summing a closed form function over the cube. That sum has $2^p$ terms, and the Markov chains used in practice, which add and remove variables, are level set sampler\index{level set sampler}s on the cube with the evidence as energy. Their mixing is expected to be good when the design is well conditioned and poor when predictors are highly collinear, which is when the posterior over models is multimodal. The vertical walk estimates the normalising constant over models by a ladder in the evidence, and the thesis says that the ladder is polynomial when the add and remove chain mixes.
